%% kapa: exam=1 round=31 year=2020 period=1 form=A subject=real_estate subject_ko=부동산학원론 seq=33 number=113
%% source: Q-Net 감정평가사 기출문제 (cst00309) artlSeq=5207219 file=extracted/1008/5207219/A형/A형/2020년도 제31회 감정평가사 1차 시험_1교시_A형.pdf
%% conversion: kordoc 4.24.51 -> tools/gen1.py (auto); review=proofread-sonnet; flags=none
\begin{question}{113}
공간으로서의 부동산에 관한 설명으로 옳지 않은 것은?
\begin{choices}
  \item 토지는 물리적 형태로서의 지표면과 함께 공중공간과 지하공간을 포함한다.
  \item 부동산활동은 3차원의 공간활동으로 농촌지역에서는 주로 지표공간이 활동의 중심이 되고, 도시지역에서는 입체공간이 활동의 중심이 된다.
  \item 지표권은 토지소유자가 지표상의 토지를 배타적으로 사용할 수 있는 권리를 말하며, 토지와 해면과의 분계는 최고만조시의 분계점을 표준으로 한다.
  \item 지중권 또는 지하권은 토지소유자가 지하공간으로부터 어떤 이익을 획득하거나 사용할 수 있는 권리를 말하며, 물을 이용할 수 있는 권리가 이에 포함된다.
  \item 공적 공중권은 일정 범위 이상의 공중공간을 공공기관이 공익목적의 실현을 위해 사용할 수 있는 권리를 말하며, 항공기 통행권이나 전파의 발착권이 이에 포함된다.
\end{choices}
\end{question}
\answer{4}
